% !TeX root = ../教材题库_全部条目_转写工作稿.tex
% Source: XJ-B1-C04-EX-P140-Y1262
% 第4章 / 4.5 函数模型及其应用 / 教材 p.140 / PDF 第147页
% OCR未校对：公式与文字均可能有误；未完成数学TeX转写；禁止直接发给学生
% 原题及相邻原解截图（核对用）：../题目原图/XJ-B1-C04-EX-P140-Y1262-01.webp
% 本文件只在转写完成后作为题干使用；OCR 条目尚不是可靠题干。
\begin{bankitem}{4.5 例题 例题2}{第4章 / 4.5 函数模型及其应用 / 教材 p.140 / PDF 第147页}{XJ-B1-C04-EX-P140-Y1262}{OCR未校对：公式与文字均可能有误；未完成数学TeX转写；禁止直接发给学生}
\begin{ocrdraft}
[教材p.140 / PDF第147页]
\par
0二0]比较3，一
\par
和3，一100工十1000在区间[0，十亠）上增长的快慢．
\par
解
\par
当>000时，
\par
有101>100十1000，因而有
\par
l.0001
\par
l.0001
\par
0．0001
\par
100工十1000101x
\par
101
\par
于是，对任意正数M>1，当>（101M）I
\par
00>1000时，有
\par
l.0001
\par
0．0001
\par
>脲101M）]
\par
100000．0001
\par
工=A亻．101x>A100十1000L
\par
这表明无论M多大，当工大到一定程度，3，一产灬就会比3，一10牖十1000的
\par
M倍还大．
\end{ocrdraft}
\end{bankitem}
